% =====================================================================
%  統計基礎 2. 推測統計
%  構成は Docs/SectionPlan.md にしたがう。
% =====================================================================
\chapter{推測統計}
\label{chap:stat-inference}

\begin{chaptergoal}
  % TODO: この章の目標
\end{chaptergoal}

% ---------------------------------------------------------------------
\section{標本分布}
\label{sec:stat-sampling-distribution}

\keywords{カイ二乗分布、標本平均と標本分散の独立性、t 分布、F 分布}

% TODO: 〔補足〕標本分布とは何か（統計量も確率変数であること）

\subsection{カイ二乗分布}
% TODO:   定義：独立な標準正規分布にしたがう変数の2乗の和
% TODO:   自由度の意味、期待値と分散、形の変化
% TODO:   ガンマ分布との関係、再生性

\subsection{標本平均と標本分散の独立性}
% TODO:   正規母集団では、標本平均 $\bar X$ と標本分散 $S^2$ が独立になる
% TODO:   $(n-1)S^2/\sigma^2$ が自由度 $n-1$ のカイ二乗分布にしたがう
% TODO:   自由度が1つ減る理由の直感（偏差の合計が必ず 0 になる）
% TODO:   証明の流れ（直交変換を使う。数学基礎「線形代数」を参照）

\subsection{t 分布}
% TODO:   定義：$Z/\sqrt{V/k}$（$Z$ は標準正規、$V$ は自由度 $k$ のカイ二乗、互いに独立）
% TODO:   母分散がわからないときの標本平均の分布として現れる
% TODO:   形の特徴（正規分布より裾が重い）、自由度を大きくすると正規分布に近づく

\subsection{F 分布}
% TODO:   定義：2つの独立なカイ二乗変数を、それぞれの自由度で割った比
% TODO:   2つの母分散の比の推測、分散分析（2.4）での使い方
% TODO:   t 分布との関係（$t^2$ は F 分布にしたがう）

% ---------------------------------------------------------------------
\section{点推定、区間推定}
\label{sec:stat-estimation}

\keywords{一致性、有効性、信頼区間と信頼係数}

\subsection{点推定の方法}
% 〔補足〕
% TODO:   モーメント法
% TODO:   最尤法：尤度関数、対数尤度、微分して 0 とおく（尤度比検定・EM 法の準備）

\subsection{推定量のよさの基準}
% TODO:   〔補足〕不偏性（復習）と平均二乗誤差（＝バイアスの2乗＋分散）
% TODO:   一致性：標本の大きさを増やすと真の値に確率収束すること。大数の法則との関係と例
% TODO:   有効性：不偏推定量の中で分散が最小であること
% TODO:     フィッシャー情報量とクラメール・ラオの下限
% TODO:     下限に達する推定量（有効推定量）の例
% TODO:   〔参考〕最尤推定量の漸近的な性質（一致性・漸近正規性）

\subsection{区間推定}
% TODO:   信頼区間と信頼係数の意味
% TODO:     「95% 信頼区間」の正しい解釈（同じ方法で何度も区間を作ると、その 95% が真の値を含む）
% TODO:     よくある誤解
% TODO:   正規母集団の母平均（母分散が既知のとき・未知のとき）
% TODO:   正規母集団の母分散（カイ二乗分布を使う）
% TODO:   母比率（正規近似）
% TODO:   2つの母平均の差、2つの母分散の比

% ---------------------------------------------------------------------
\section{汎用的な検定}
\label{sec:stat-general-tests}

\keywords{尤度比検定、ノンパラメトリック検定、ウィルコクソン検定、並べ替え検定}

\subsection{尤度比検定}
% TODO:   尤度比の定義と考え方（帰無仮説のもとでの最大尤度と、全体での最大尤度の比）
% TODO:   〔参考〕ネイマン・ピアソンの補題（最も検出力の高い検定）
% TODO:   $-2\log(\text{尤度比})$ が近似的にカイ二乗分布にしたがうこと（ウィルクスの定理）
% TODO:   例：正規母集団の母平均の検定、二項分布の比率の検定

\subsection{ノンパラメトリック検定}
% TODO:   パラメトリック検定との違い（母集団の分布を仮定しない）
% TODO:   長所（外れ値に強い）と短所（検出力が下がることがある）
% TODO:   順位を使うという考え方、符号検定

\subsection{ウィルコクソン検定}
% TODO:   順位和検定：対応のない2群の比較（マン・ホイットニーの U 検定と同じもの）
% TODO:   符号付き順位検定：対応のある2群の比較
% TODO:   小さなデータでの計算例、標本が大きいときの正規近似

\subsection{並べ替え検定}
% TODO:   考え方：帰無仮説が正しければ、群のラベルを入れかえても分布は変わらない
% TODO:   手順と、正確な p 値の求め方
% TODO:   組合せの数が多いときはランダムに入れかえて近似する（計算統計 4 章への橋渡し）

% ---------------------------------------------------------------------
\section{種々の検定}
\label{sec:stat-various-tests}

\keywords{一元配置分散分析、二元配置分散分析、交互作用、適合度検定}

\subsection{一元配置分散分析}
% TODO:   モデル：データ ＝ 全体平均 ＋ 群の効果 ＋ 誤差
% TODO:   平方和の分解（全体の平方和 ＝ 群間平方和 ＋ 群内平方和）
% TODO:   自由度、分散分析表、F 検定
% TODO:   t 検定との関係（2群のとき）

\subsection{二元配置分散分析}
% TODO:   2つの要因の主効果
% TODO:   繰り返しのない場合と、繰り返しのある場合
% TODO:   平方和の分解と分散分析表

\subsection{交互作用}
% TODO:   意味：一方の要因の効果が、もう一方の要因の水準によって変わること
% TODO:   交互作用プロットによる見方
% TODO:   繰り返しのある二元配置での交互作用の検定

\subsection{適合度検定}
% TODO:   ピアソンのカイ二乗統計量
% TODO:   自由度の決め方（推定したパラメータの数だけ減る）
% TODO:   例：サイコロの公平性、ポアソン分布への当てはまり
% TODO:   分割表の独立性の検定
% TODO:   〔参考〕尤度比検定（G 検定）との関係

% ---------------------------------------------------------------------
\section{多重比較}
\label{sec:stat-multiple-comparison}

\keywords{ボンフェロニ補正}

% TODO: 多重性の問題：検定を何回もくり返すと、どこかで誤って棄却する確率が大きくなる（$1-(1-\alpha)^m$ の計算例）
% TODO: ファミリーワイズ・エラー率（FWER）

\subsection{ボンフェロニ補正}
% TODO:   方法：有意水準を検定の回数 $m$ で割る（$\alpha/m$）
% TODO:   根拠：確率の和の不等式（ブールの不等式）
% TODO:   欠点：保守的（棄却しにくくなる）
% TODO: 〔参考〕ホルム法、テューキー法、偽発見率（FDR）

\begin{chaptersummary}
  % TODO: この章のまとめ
\end{chaptersummary}
